Программа должна реализовывать генерацию таблиц набора данных TPC-H с учетом детерминированных правил синтеза значений, зависимостей между таблицами и возможности независимой обработки партиций. При необходимости, если этого будет требовать реализация сценария хранения, программа должна обеспечивать формирование промежуточного колонoчного представления данных в памяти и запись результата в поддерживаемые форматы хранения с различными параметрами сериализации.

Программа также должна реализовывать возможность выбора стратегии записи, позволяющую пользователю либо максимизировать пропускную способность за счет записи партиций в отдельные файлы, либо формировать один итоговый файл таблицы при сохранении параллельной генерации, что позволяет использовать программу в различных сценариях тестирования и исследования.

Разрабатываемая программа должна: 

\begin{enumerate}
    \item генерировать и сохранять таблицы набора данных TPC-H;
    \item поддерживать генерацию следующих таблиц:
    \begin{enumerate}
        \item region;
        \item nation;
        \item supplier;
        \item part;
        \item partsupp;
        \item customer;
        \item orders;
        \item lineitem.
    \end{enumerate}
    \item изменять масштаб генерируемого набора данных посредством параметра scale factor;
    \item формировать данные батчами фиксированного размера в колонoчном представлении;
    \item позволять пользователю выбирать набор генерируемых таблиц;
    \item переводить сгенерированные данные в различные форматы хранения:
    \begin{enumerate}
        \item parquet;
        \item orc;
        \item lance;
        \item vortex.
    \end{enumerate}
    \item поддерживать различные стратегии записи:
    \begin{enumerate}
        \item parallel-files;
        \item single-file-ordered.
    \end{enumerate}
    \item позволять пользователю управлять параметрами генерации и записи:
    \begin{enumerate}
        \item размером батча;
        \item количеством рабочих потоков;
        \item емкостью очереди записи;
        \item параметрами компрессии и партиционирования.
    \end{enumerate}
    \item обеспечивать детерминированность результата при изменении конфигурации партиционирования и числа потоков;
    \item выводить список доступных таблиц и форматов;
    \item выполнять проверку корректности параметров запуска и выдавать диагностические сообщения пользователю.
\end{enumerate}
